% concatenated from Source.tex
\documentclass[12pt]{amsart}
\usepackage[T1]{fontenc}
\usepackage[a4paper,margin=1in]{geometry}
\usepackage{graphicx}
\usepackage{amssymb}
\usepackage{stmaryrd}
\usepackage{enumitem}
\DeclareGraphicsRule{.tif}{png}{.png}{`convert #1 `dirname
#1`/`basename #1 .tif`.png}


\usepackage{tikz}
\usepackage{float}
\usepackage{algorithm}
%\usepackage{algorithmic}
\usepackage{algpseudocode}
\usepackage{url}
\usepackage{mathtools}
\usepackage{amsmath,amsfonts,amsthm}
\usepackage{mathrsfs} 
\newcommand{\case}[1]{\paragraph*{Case #1:}}
%\usepackage[pdfencoding=auto,unicode,psdextra]{hyperref} 

\usepackage{amsthm}

\theoremstyle{plain}
\newtheorem{theorem}{Theorem}[section]
\newtheorem{lemma}[theorem]{Lemma}
\newtheorem{proposition}[theorem]{Proposition}
\newtheorem{corollary}[theorem]{Corollary}
\newtheorem{conjecture}[theorem]{Conjecture}

\theoremstyle{definition}
\newtheorem{definition}{Definition}
\newtheorem{example}{Example}
\newtheorem{question}{Question}
\newtheorem{openproblem}{Open Problem}

\theoremstyle{remark}
\newtheorem{remark}[theorem]{Remark}
%================
\usepackage{multirow}

\newcommand{\fphi}{\varphi}
\newcommand{\CE}{\operatorname{CE}}
\newcommand{\ice}{\operatorname{ice}}
\newcommand{\aice}{\operatorname{ice}^*}
\newcommand{\N}{\mathbb{N}}
\newcommand{\Z}{\mathbb{Z}}
\newcommand{\AlphA}{\mathcal{A}}
\newcommand{\FD}{\mathfrak{F}}
\renewcommand{\qedsymbol}{$\blacksquare$}
\newcommand{\doi}[1]{\href{https://doi.org/#1}{DOI: #1}}

\usepackage[colorlinks,
linkcolor=blue,
anchorcolor=blue,
citecolor=blue, 
pdfencoding=auto,
unicode,psdextra
]{hyperref}

\author{J. Hamoud}
\address{\textbf{Jasem Hamoud:} Department of Discrete Mathematics, Moscow Institute of Physics and Technology}
\email{hamoud.math@gmail.com}
\thanks{ }

\title[Duality in Biperiodic Fibonacci Words]{Duality in Biperiodic Fibonacci Words Substitution Frequencies and Combinatorial Invariants}

\date{}
\begin{document}


\begin{abstract}
In this paper, a natural duality on the family of biperiodic Fibonacci words $\mathfrak{F}^{(a,b)}$ generated by the directive sequence $(a,b,a,b,\dots)$. Both of $\mathfrak{F}^{(a,b)}$ and $\mathfrak{F}^{(b,a)}$ are related by the explicit morphism $\sigma_a:0\mapsto 0^a1$, $1\mapsto 0$, establishing a precise substitutional correspondence between them.
We compute the exact letter frequencies, give a complete description of the return words for each letter, prove the existence of arbitrarily long palindromic prefixes, and determine the continued fraction expansion of the slope $\theta^{(a,b)}$. 

This findings reveal that the apparent asymmetry in several invariants  arises uniformly from the length-redistribution mechanism induced by the morphism $\sigma_a$. 
\end{abstract}

\maketitle


\noindent\rule{15.9cm}{1.0pt}

\noindent
\textbf{
Keywords:} Fibonacci word; Sturmian word; biperiodic Fibonacci word; standard sequence; continued fraction; critical exponent; return words; palindromic complexity; combinatorics on words.


\medskip

\noindent
{\bf
MSC 2020:} 68R15 (Primary); 11B39, 37B10, 11A55 (Secondary).

\medskip

\noindent
{\bf
UDC:} 519.1



\noindent\rule{15.9cm}{1.0pt}

%======================
\section{Introduction}~\label{sec01intro}
Let $\AlphA=\{0,1\}$ be a binary alphabet. A \emph{finite word} over
$\AlphA$ is an element of the free monoid $\AlphA^{*}$, and an
\emph{infinite word} is an element of $\AlphA^{\N}$. Among all infinite binary words, the
\emph{Fibonacci word}
\[
  \FD = 0100101001001010010100100101001001\cdots,
\]
obtained  of the morphism $\sigma\colon 0\mapsto 01,\;
1\mapsto 0$.  For a finite word
$w=w_0w_1\cdots w_{k-1}$ denote by $|w|=k$ for its length, and $|w|_c$
for the number of occurrences of the letter $c\in\AlphA$ in $w$, where
$|w|=|w|_0+|w|_1$. A finite word $u$ is a \emph{factor} of a (finite or
infinite) word $w$ if $u=w_iw_{i+1}\cdots w_{i+|u|-1}$ with $i\geqslant 0$;
it is a \emph{prefix} if $i=0$. Denote by $\mathcal{L}(w)$ for the set of
all factors of $w$ (its \emph{language}), and $\mathcal{L}_n(w)$ for the
factors of length $n$. The \emph{factor complexity function} of $w$ is
$p_w(n)=|\mathcal{L}_n(w)|$.
A finite or infinite word $w$ is a \emph{palindrome} if it is invariant
under reversal; we write $\widetilde{w}$ for the reversal of a finite
word $w$.
 

A \emph{morphism} $\tau\colon\AlphA^{*}\to\AlphA^{*}$ is a monoid
homomorphism, determined by the finite data $\tau(0),\tau(1)$, and
extended to infinite words coordinatewise in the obvious way whenever
the extension is well defined. Following Ram\'irez, and Rubiano~\cite{Ramirez2015}, we fix integer parameters $a,b\ge1$ and seeds
\[
  \FD^{(a,b)}_0 = 1, \qquad \FD^{(a,b)}_1 = 0,
\]
and define $\FD^{(a,b)}_n$ for $n\geqslant 2$ by the recurrence already stated
in~\eqref{eqq01biperiodicrecurrence}. The
\emph{biperiodic word} is generated by the two parameter recurrence
\begin{equation}~\label{eqq01biperiodicrecurrence}
  \FD^{(a,b)}_n =
  \begin{cases}
    \bigl(\FD^{(a,b)}_{n-1}\bigr)^a\, \FD^{(a,b)}_{n-2} & \text{if } n \text{ is even},\\[4pt]
    \bigl(\FD^{(a,b)}_{n-1}\bigr)^b\, \FD^{(a,b)}_{n-2} & \text{if } n \text{ is odd},
  \end{cases}
\end{equation}
with $a,b\geqslant 1$ integers and suitable seeds $\FD^{(a,b)}_0,\FD^{(a,b)}_1$.


D.~Abdullah, and M.~Meisami~\cite{Abdullah2023Meisami} had  established that the factor complexity of the infinite word $\mathfrak{F}_b$ is full and the arithmetic complexity of the infinite word $\mathfrak{F}_b$ is full (see~\cite{Abdullah2025Hamoud}). 
The smallest critical exponent attained by any Sturmian word~\cite{Cassaigne2008, DamanikLenz2002, MousaviShallit2015} where its \emph{critical exponent} $\CE(\FD) = 2+\fphi$.
 The least repetitive of all aperiodic balanced binary
words, which is what makes it the canonical hard instance for
repetition-detection algorithms~\cite{AllouecheShallit2003} and the
extremal case in Dejean type repetition threshold
problems~\cite{CurrieRampersad2008}.

\medskip

The $k$-bonacci words, the Tribonacci
word, and the family relevant to this paper, the \emph{biperiodic
Fibonacci words} $\FD^{(a,b)}$ of Ram\'irez, and Rubiano~\cite{Ramirez2015} (see also~\cite{EdsonYayenie2009, MignosiPirillo1992})had showed that the letter frequencies of
$\FD^{(a,b)}$ converge to
\begin{equation}~\label{eqq1biperiodicfreqs}
  \lim_{n\to\infty}\frac{|\FD^{(a,b)}_n|_1}{|\FD^{(a,b)}_n|} = \frac{1}{1+\AlphA a},
  \qquad
  \lim_{n\to\infty}\frac{|f^{(a,b)}_n|_0}{|f^{(a,b)}_n|} = \frac{\AlphA a}{1+\AlphA a},
\end{equation}
where
\begin{equation}~\label{eqq1AlphAadef}
  \AlphA a = \AlphA a(a,b) = \frac{\sqrt{a^2b^2+4ab}+ab}{2},
\end{equation}
and that the critical exponent of $\FD^{(a,b)}$ is
\begin{equation}~\label{eqq1biperiodiccritical}
  \CE\bigl(\FD^{(a,b)}\bigr) = 2+\frac{\AlphA a}{a}+\frac{b}{\AlphA a}.
\end{equation}
Actually, if $a=b=1$ in~\eqref{eqq1AlphAadef} gives $\AlphA a=\fphi$. Then, the relation~\eqref{eqq1biperiodicfreqs} recovers the classical frequencies
$1/\fphi$ and $1/\fphi^2$, and the relation~\eqref{eqq1biperiodiccritical} recovers
$\CE(f)=2+\fphi$.  
According to~\eqref{eqq1AlphAadef}, it establishes that
$\AlphA a(a,b)=\AlphA a(b,a)$. Then,  $\FD^{(a,b)}$ is \emph{symmetric} under exchanging
the two periods $a$ and $b$ where the words $\FD^{(a,b)}$ and $\FD^{(b,a)}$ have
identical asymptotic letter frequencies for every choice of $a,b$. 
Based on the critical exponent relation~\eqref{eqq1biperiodiccritical} gives
\begin{equation}~\label{eqq1cedifference}
  \CE(a,b) - \CE(b,a)
  = (b-a)\left(\frac{\AlphA a}{ab} + \frac{1}{\AlphA a}\right),
\end{equation}
which is strictly positive whenever $b>a$ and strictly negative whenever
$b<a$, and vanishes only when $a=b$. 


The asymmetry in~\eqref{eqq1cedifference} is not an
artefact of the specific closed form~\eqref{eqq1biperiodiccritical}, but
a shadow of a precise structural relationship between $\FD^{(a,b)}$ and
$\FD^{(b,a)}$ at the level of the words themselves. Unfolding the
recurrence~\eqref{eqq01biperiodicrecurrence} by one step shows that the
word $\FD^{(a,b)}_{n+1}$, viewed as a sequence indexed from level $n+1$
rather than level $n$, obeys the \emph{same} recurrence with the roles
of $a$ and $b$ interchanged. This observation by considering elementary at the level
of the recurrence, but with nontrivial consequences once the initial
data and the necessary letter relabelling are tracked carefully through
the limit $n\to\infty$, suggests the existence of an explicit
correspondence $\theta$, described precisely in
Section~\ref{sec01continuedfraction}, such that
\[
 \FD^{(b,a)}=\theta\bigl(\FD^{(a,b)}\bigr)
\]
as infinite words, up to a shift and a bounded prefix correction. We
call this the \emph{Parity-Shift Duality} of the biperiodic Fibonacci
family. 

A subsequent study extended this geometric side of the theory with a
self-similar decomposition and a family of ``biperiodic Fibonacci
snowflake'' polyominoes~\cite{FibSnowflake2024, Dumaine2009}. 
Critical exponents and their initial asymptotic refinements~\cite{Cassaigne2008, MousaviShallit2015}, balance and abelian
complexity for Sturmian and other morphic
words~\cite{BalkovaBrindaTurek2011, RichommeSaariZamboni2010}, and the
$S$-adic continued fraction coding of standard
sequences~\cite{AllouecheShallit2003, Lothaire2002}. The duality relationship
between $\FD^{(a,b)}$ and $\FD^{(b,a)}$ that organises our results appears
not to have been observed.


Section~\ref{sec02prelim} fixes notation and recalls the definitions of
critical exponent (global, initial, and asymptotic), balance, abelian
complexity, and palindromic complexity. Section~\ref{sec01sadic}
constructs the $S$-adic representation of $\FD^{(a,b)}$ via the periodic
continued fraction $[a,b,a,b,\ldots]$. Section~\ref{sec04ceduality} states
and proves the Parity-Shift Duality Theorem, subsection~\ref{sec5duality}
derives the duality for the four repetition exponent
variants and recovers combinatorially. 
Section~\ref{sec01continuedfraction} determine the continued fraction expansion of the
slope of the biperiodic Fibonacci word $\FD^{(a,b)}$, i.e.\ of the
irrational number.  
%======================
\section{Preliminaries}~\label{sec02prelim}

Since $\FD^{(a,b)}_n$ is a prefix of $\FD^{(a,b)}_{n+1}$ for every $n$, the sequence $\bigl(\FD^{(a,b)}_n\bigr)_{n\ge0}$
converges to a well-defined infinite word
\[
  \FD^{(a,b)} = \lim_{n\to\infty} \FD^{(a,b)}_n \in \AlphA^{\N},
\]
the \emph{biperiodic Fibonacci word}. Assume that $\sigma_a\colon 0\mapsto
0^a1,\,1\mapsto 0$ and $\sigma_b\colon 0\mapsto 0^b1,\,1\mapsto 0$ for
the two morphisms whose alternating application generates
$\FD^{(a,b)}$; this alternation is made precise via the $S$-adic
formalism in Section~\ref{sec01sadic}.  
For a finite word $w\ne\varepsilon$, the \emph{exponent} of $w$ is
$e(w) = |w|/\pi(w)$, where $\pi(w)$ is the smallest period of $w$. A
finite word $y$ is an \emph{$\AlphA a$-power} of a primitive word $x$, for
rational $\AlphA a>1$, if $y=x^{\lfloor\AlphA a\rfloor}x_p$ where $x_p$ is a
prefix of $x$ and $|y|=\AlphA a|x|$. 

\begin{definition}~\label{def01ce}
The critical exponent of $w$ is
\[
  \CE(w) = \sup\{\, e(u) : u \text{ a factor of } w \,\} \in (1,\infty].
\]
\end{definition}
Based on Definition~\ref{def01ce}, the initial critical exponent of $w$ is the supremum of $e(u)$
over \emph{prefixes} $u$ of $w$, $\ice(w) = \sup\{\, e(u) : u \text{ a prefix of } w\,\},$
and the asymptotic critical exponent of $w$ restricts the
supremum defining $\CE(w)$ to factors of unbounded length, $\aice(w) = \lim_{n\to\infty}\ \sup\{\, e(u) : u \in \mathcal{L}(w),\ |u|\geqslant n \,\}.$
\begin{definition}~\label{def01balance}
An infinite word $w$ is $C$-balanced ($C\in\N$) if for every pair
of factors $u,v\in\mathcal{L}(w)$ with $|u|=|v|$, we have
$\bigl||u|_1-|v|_1\bigr|\le C$. The \emph{balance function} $B(n)$
records the smallest $C$ for which this holds among all factor pairs of
length $n$, and $w$ is \emph{balanced} if $\sup_n B(n)<\infty$.
\end{definition}

Sturmian words are exactly the aperiodic $1$-balanced binary
words~\cite{MorseHedlund1940}; since $\FD^{(a,b)}$ is generated by an
alternation of two \emph{different} morphisms rather than the iteration
of one, it need not be Sturmian for $a\ne b$, and its balance function
is not known in closed form for the biperiodic family.

\begin{definition}~\label{def01abelian}
Two factors $u,v\in\mathcal{L}_n(w)$ are abelian equivalent,
$u\sim_{ab}v$, if $|u|_c=|v|_c$ for every letter $c\in\AlphA$. The
\emph{abelian complexity function} is
$\mathcal{AC}_w(n)=|\mathcal{L}_n(w)/\!\sim_{ab}|$.
\end{definition}

Abelian complexity is known explicitly for Sturmian
words~\cite{RichommeSaariZamboni2010} and for the words associated with
quadratic Parry numbers~\cite{BalkovaBrindaTurek2011}, but has not been
computed for $\FD^{(a,b)}$ as an explicit function of $(a,b)$.

\begin{definition}~\label{def01palcomplexity}
The palindromic complexity function of $w$ is
$P_w(n)=|\{u\in\mathcal{L}_n(w): u=\widetilde u\}|$, the number of
distinct length-$n$ palindromic factors.
\end{definition}

A certain derived word
$\Phi(\FD^{(a,b)}_n)$ is a palindrome for every $n$ and all $a\geqslant 2,\,
b\geqslant 1$ (see~\cite{Ramirez2015}). In fact, this is a statement about one distinguished
factor per level $n$, not about $P_{\FD^{(a,b)}}(n)$ itself, which remains
uncomputed for the family.

\begin{definition}~\label{def01central}
For a word $w$ with $|w|\ge2$, write $w^{-}$ for $w$ with its last two
letters deleted.
\end{definition}

\begin{theorem}[Vuillon~\cite{Vuillon2001}]
\label{thm01vuillon}
An aperiodic binary word $w$ is Sturmian if and only if every
$u\in\mathcal{L}(w)$ has exactly two return words in $w$.
\end{theorem}

\subsection{Problem Statement}~\label{subsec00problem}

Two invariants of $\FD^{(a,b)}$ built from
the \emph{same} algebraic quantity $\AlphA a(a,b)$ behave differently
under the exchange $a\leftrightarrow b$.


\begin{conjecture}~\label{conj01dualityinformal}
There is an explicit, effectively computable map $\theta$ on infinite
binary words, built from a letter relabelling and a bounded prefix
correction, such that
\[
  \FD^{(b,a)} = \theta\bigl(\FD^{(a,b)}\bigr)
  \qquad \text{for all } a,b\geqslant 1.
\]
\end{conjecture}

This yields the four questions that structure the remainder of the paper.

\begin{question}~\label{q01ce}
Does the duality of Conjecture~\ref{conj01dualityinformal} extend to
the initial critical exponent $\ice$ and the asymptotic critical
exponent $\aice$? If so, what are $\ice(\FD^{(a,b)})$ and
$\aice(\FD^{(a,b)})$ in closed form; if not, where does duality break down?
\end{question}

\begin{question}~\label{q02balance}
What is the balance function $B(n)$ of $\FD^{(a,b)}$, and is $B$ invariant
under $\theta$? In particular, for which $(a,b)$ is $\FD^{(a,b)}$ balanced
in the sense of Definition~\ref{def01balance}, and does the answer
depend on $\min(a,b)$, $\max(a,b)$, or only on $\AlphA a(a,b)$?
\end{question}

\begin{question}~\label{q03abelian}
What is the abelian complexity function $\mathcal{AC}_{\FD^{(a,b)}}(n)$ as
an explicit function of $a,b,n$, and does $\theta$ intertwine
$\mathcal{AC}_{\FD^{(a,b)}}$ with $\mathcal{AC}_{\FD^{(b,a)}}$?
\end{question}

\begin{question}~\label{q04palindrome}
What is the palindromic complexity function $P_{\FD^{(a,b)}}(n)$, how does
it relate to the known palindromicity of $\Phi(\FD^{(a,b)}_n)$, and is it
$\theta$-invariant?
\end{question}

%========== SECTION 3 =================================== 
%========================================================
\section{The S-adic Continued Fraction Model}~\label{sec01sadic}

The recurrence~\eqref{eqq01biperiodicrecurrence} defining $\FD^{(a,b)}$ has
exactly the shape of the classical \emph{singular sequence}
construction used to build Sturmian words from continued fractions. 
For an integer $d\ge1$, let $\sigma_d\colon\AlphA ^{*}\to\AlphA ^{*}$ be the
\emph{elementary morphism}
\[
  \sigma_d(0) = 0^d1, \qquad \sigma_d(1) = 0.
\]
Hence, each $\sigma_d$ is non-erasing and injective on $\AlphA ^{\N}$, and it extends continuously to a map $\AlphA ^{\N}\to\AlphA ^{\N}$: if $u_n\to u$
in the prefix topology, then $\sigma_d(u_n)\to\sigma_d(u)$.

\begin{definition}~\label{def01standardsequence}
Given a sequence of positive integers $(d_k)_{k\ge1}$, the associated \emph{standard sequence} $(s_k)_{k\ge-1}$ is
defined by
\[
  s_{-1} = 1, \qquad s_0 = 0, \qquad
  s_k = s_{k-1}^{\,d_k}\,s_{k-2} \quad (k\ge1).
\]
\end{definition}

Since $s_{k-2}$ is a prefix of $s_{k-1}$ for every $k$, $s_k$ is a
prefix of $s_{k+1}$ for every $k\ge-1$, and $(s_k)$ converges to a
well-defined infinite word
\[
  s^{(d_1,d_2,d_3,\ldots)} = \lim_{k\to\infty} s_k \in \AlphA ^{\N}.
\]
It is a classical fact  that $s^{(d_1,d_2,\ldots)}$ is a \emph{standard Sturmian word} for
\emph{every} choice of the directive sequence $(d_k)_{k\ge1}$, and that
every standard Sturmian word arises this way; the slope of
$s^{(d_1,d_2,\ldots)}$ is the irrational number with continued fraction
expansion built from $(d_k)$. In particular periodicity of $(d_k)$ is
not required for $s^{(d_1,d_2,\ldots)}$ to be Sturmian: it is Sturmian
for \emph{any} sequence of positive integers, and merely happens to have
\emph{quadratic irrational} slope precisely when $(d_k)$ is eventually
periodic.

A second classical fact we will use is the \emph{desubstitution
identity}: for every directive sequence,
\begin{equation}  \label{eq01desubgeneral}
  s^{(d_1,d_2,d_3,\ldots)} = \sigma_{d_1}\bigl(s^{(d_2,d_3,\ldots)}\bigr).
\end{equation}

\begin{proposition}~\label{prop01fabisstandard}
For every $a,b\geqslant 1$, the biperiodic Fibonacci word $f^{(a,b)}$
coincides with the standard sequence $s^{(d_1,d_2,d_3,\ldots)}$ with
directive sequence
\[
  (d_1,d_2,d_3,d_4,\ldots) = (a,b,a,b,\ldots), \qquad
  d_k = \begin{cases} a & k \text{ odd},\\ b & k \text{ even}.\end{cases}
\]
\end{proposition}
\begin{proof}
According to Definition~\ref{def01standardsequence} and the seeds
$\FD^{(a,b)}_0=1=s_{-1}$, $\FD^{(a,b)}_1=0=s_0$, and the
recurrence~\eqref{eqq01biperiodicrecurrence} reindexed as $k=n-1$,  for
$n\geqslant 2$,  we obtain $\FD^{(a,b)}_n = s_{n-1}$ requires $s_k = s_{k-1}^{d_k}s_{k-2}$
with $d_k=a$ when $n=k+1$ is even. An immediate
induction on $n$ then gives $\FD^{(a,b)}_n=s_{n-1}$ for all $n\geqslant 0$, and
by considering  $n\to\infty$ gives $\FD^{(a,b)}=s^{(a,b,a,b,\ldots)}$.
\end{proof}

Corollary~\ref{cor01sturmian} had established $1$-balanced in the sense of
Definition~\ref{def01balance}, and its slope is a quadratic irrational.

\begin{corollary}~\label{cor01sturmian}
For every $a,b\geqslant 1$, $\FD^{(a,b)}$ is a standard Sturmian word.
\end{corollary}

This has a direct
consequence for Question~\ref{q02balance}: since every aperiodic
Sturmian word is exactly $1$-balanced and has abelian complexity
identically $2$~\cite{RichommeSaariZamboni2010}, both $B(n)\equiv1$ and
$\mathcal{AC}_{\FD^{(a,b)}}(n)\equiv2$ hold for all $a,b\geqslant 1$, and
duality under $a\leftrightarrow b$ for these two invariants is trivial.

Definition~\ref{def01standardsequence} attaches to any directive
sequence $(d_k)$ the continued fraction $$[d_1,d_2,d_3,\ldots] :=
d_1+\cfrac{1}{d_2+\cfrac{1}{d_3+\cdots}}.$$
For the alternating sequence
of Proposition~\ref{prop01fabisstandard} continued fraction is
purely periodic with period $(a,b)$, and we denote its value by $\chi(a,b) := [a,b,a,b,\ldots].$

\begin{proposition}~\label{prop01AlphAacf}
For all $a,b\geqslant 1$, $\AlphA a(a,b) = b\cdot\chi(a,b).$
\end{proposition}
\begin{proof}
Assume that $x=\chi(a,b)$, the periodicity $[a,b,a,b,\ldots] =
a+1/[b,a,b,a,\ldots]$ together with $[b,a,b,a,\ldots] = b+1/x$ satisfies
\[
  x = a+\frac{1}{\,b+\frac1x\,} = a + \frac{x}{bx+1} = \frac{abx+a+x}{bx+1}.
\]
Since $x(bx+1) = abx+a+x$, where $bx^2-abx-a=0,$
whose positive root is $x = (ab+\sqrt{a^2b^2+4ab})/2b.$ Thus, by multiplying through by $b$ gives $bx =(ab+\sqrt{a^2b^2+4ab})/2,$ 
which is exactly $\AlphA a(a,b)$ as given in~\eqref{eqq1AlphAadef}.
\end{proof}


Proposition~\ref{prop01AlphAacf} identifies $\AlphA a(a,b)$, and the
letter frequencies~\eqref{eqq1biperiodicfreqs}, as a value of the
periodic continued fraction attached to $\FD^{(a,b)}$ by considering
Proposition~\ref{prop01fabisstandard}. 
Assume that $x=a+1/y$ with $y=b+1/x$, is precisely a statement about the
two shifts of one periodic continued fraction, and it is this
shift relation  at the level of directive sequences, by considering the
duality between $\FD^{(a,b)}$ and $\FD^{(b,a)}$. Thus, $\FD^{(a,b)}$ has directive
sequence $(a,b,a,b,\ldots)$ and $\FD^{(b,a)}$ has directive sequence
$(b,a,b,a,\ldots)$ where $\FD^{(b,a)}$ is exactly $\FD^{(b,a)}$ with its leading
term omitted and emphasized that $\FD^{(a,b)} = \sigma_a\bigl(\FD^{(b,a)}\bigr)$. 

\begin{theorem}~\label{thm01desub}
For all $a,b\geqslant 1$ and all $n\geqslant 0$,
\[
  \sigma_a\bigl(\FD^{(b,a)}_n\bigr) = \FD^{(a,b)}_{n+1},
\]
if and only if  $\sigma_a\bigl(\FD^{(b,a)}\bigr) = \FD^{(a,b)}$ for large $n$.
\end{theorem}
\begin{proof}
By using induction mathematics, for $n=0,1$, we have $\FD^{(b,a)}_0=1$ and $\sigma_a(1)=0=\FD^{(a,b)}_1$. Then, $\FD^{(b,a)}_1=0$ and $\sigma_a(0)=0^a1=\FD^{(a,b)}_1{}^{a}\FD^{(a,b)}_0$. Thus, for $n=2$, $\FD^{(a,b)}_2=\bigl(\FD^{(a,b)}_1\bigr)^a \FD^{(a,b)}_0$, and $\FD^{(a,b)}_2=0^a1=\sigma_a\bigl(\FD^{(b,a)}_1\bigr)$.

Now, let $n\geqslant 2$ and suppose that $\sigma_a(\FD^{(b,a)}_{n-1})=\FD^{(a,b)}_n$ and $\sigma_a(\FD^{(b,a)}_{n-2})=\FD^{(a,b)}_{n-1}$. Since $\FD^{(b,a)}$ is the word obtained from~\eqref{eqq01biperiodicrecurrence} with the terms of $a$ and $b$ exchanged as
\[
  \FD^{(b,a)}_n =
  \begin{cases}
    \bigl(\FD^{(b,a)}_{n-1}\bigr)^{b}\, \FD^{(b,a)}_{n-2} & n \text{ even},\\[2pt]
    \bigl(\FD^{(b,a)}_{n-1}\bigr)^{a}\, \FD^{(b,a)}_{n-2} & n \text{ odd}.
  \end{cases}
\]
Thus, according to the morphism $\sigma_a$  and the inductive hypothesis, it follows that 
\[
  \sigma_a\bigl(\FD^{(b,a)}_n\bigr) =
  \begin{cases}
    \bigl(\FD^{(a,b)}_{n}\bigr)^{b}\, \FD^{(a,b)}_{n-1} & n \text{ even},\\[2pt]
    \bigl(\FD^{(a,b)}_{n}\bigr)^{a}\, \FD^{(a,b)}_{n-1} & n \text{ odd}.
  \end{cases}
\]
On the other hand, by applying~\eqref{eqq01biperiodicrecurrence} to
$\FD^{(a,b)}_{n+1}$ satisfies
\begin{enumerate}
    \item If $n$ is even then $n+1$ is odd, so $\FD^{(a,b)}_{n+1}=\bigl(\FD^{(a,b)}_n\bigr)^b \FD^{(a,b)}_{n-1}$;
    \item If $n$ is odd then $n+1$ is even, so $\FD^{(a,b)}_{n+1}=\bigl(\FD^{(a,b)}_n\bigr)^a
\FD^{(a,b)}_{n-1}$.
\end{enumerate}
 These are exactly the two cases above, so $\sigma_a(\FD^{(b,a)}_n)=\FD^{(a,b)}_{n+1}$, completing the induction.

Since $\sigma_a$ is non-erasing, it is continuous with respect to the prefix topology on $\AlphA ^{\N}$. Then, if $u_n\to u$ then $\sigma_a(u_n)\to\sigma_a(u)$. As $\FD^{(b,a)}_n\to
\FD^{(b,a)}$, we conclude $\sigma_a(\FD^{(b,a)})=\lim_n
\sigma_a(\FD^{(b,a)}_n)=\lim_n \FD^{(a,b)}_{n+1}=\FD^{(a,b)}$.
\end{proof}

Conjecture~\ref{conj01dualityinformal} holds in the following sharp,
correction-free form: the map $\theta=\sigma_a$ satisfies $\FD^{(b,a)} \ \longmapsto\ \sigma_a\bigl(\FD^{(b,a)}\bigr)=\FD^{(a,b)}$
exactly, with no letter relabelling and no bounded prefix correction
needed. Symmetrically, $\FD^{(a,b)}$ determines $\FD^{(b,a)}$ as the unique
infinite word $w$ with $\sigma_a(w)=\FD^{(a,b)}$. 


%==========================comlete steps ==================
%==========================================================
\section{Frequency Duality and Critical Exponents}~\label{sec04ceduality}

This section explores the duality induced by the substitution $\sigma_a$ (resp.\ $\sigma_b$) between the two families of infinite words $\FD^{(a,b)}$ and $\FD^{(b,a)}$. In particular, we apply Theorem~\ref{thm01desub} to relate the letter frequencies and critical exponents of these words, thereby establishing a precise correspondence between their combinatorial and statistical properties.

\begin{lemma}~\label{lem01length}
For every finite word $x\in\AlphA^{*}$ and every $d\geqslant 1$, $|\sigma_d(x)| = d\,|x|_0 + |x|.$
\end{lemma}
\begin{proof}
Let $x$ be a finite word in $\AlphA^{*}$ and let $d\geqslant 1$. Since $\sigma_d(0)=0^d1$ has length $d+1$ and $\sigma_d(1)=0$ has length $1$. Then, since $\sigma_d$ is a morphism, $|\sigma_d(x)| = (d+1)|x|_0 + 1\cdot|x|_1 =
d|x|_0 + (|x|_0+|x|_1) = d|x|_0+|x|$.
\end{proof}

\begin{lemma}~\label{lem01primitivity}
If $u\in\AlphA^{*}$ is primitive,
then $\sigma_d(u)$ is primitive for every $d\geqslant 1$.
\end{lemma}
\begin{proof}
This is a standard consequence of the recognizability of primitive
substitutions applied to the
non-erasing, injective morphism $\sigma_d$ (see
Lothaire~\cite{Lothaire2002}, Ch.~2). 
\end{proof}

\begin{theorem}~\label{thm02freqduality}
For all $a,b\geqslant 1$, the letter frequencies of $f^{(a,b)}$ are
\begin{equation}~\label{eq01freqcorrected}
  \operatorname{freq}_1\bigl(\FD^{(a,b)}\bigr)=\frac{b}{\alpha+b},
  \qquad
  \operatorname{freq}_0\bigl(\FD^{(a,b)}\bigr)=\frac{\alpha}{\alpha+b}.
\end{equation}
\end{theorem}
\begin{proof}
Assume that  $x=\FD^{(b,a)}_{n-1}$. According to Theorem~\ref{thm01desub},
$\FD^{(a,b)}_n=\sigma_a(x)$. Since $\sigma_a(0)=0^a1$ contributes one
letter $1$ and $\sigma_a(1)=0$ contributes none, 
\[
  |\FD^{(a,b)}_n|_1=|x|_0, \qquad
  |\FD^{(a,b)}_n|_0=a|x|_0+|x|_1,
\]
according to Lemma~\ref{lem01length}, $|\FD^{(a,b)}_n|=a|x|_0+|x|=(a+1)|x|_0+|x|_1$. Let $P=\operatorname{freq}_0(\FD^{(b,a)})$ and
$Q=\operatorname{freq}_1(\FD^{(b,a)})=1-P$ denote the (a priori unknown)
limiting frequencies of $\FD^{(b,a)}$. Then, $|x|=|\FD^{(b,a)}_{n-1}|$ and by considering $n\to\infty$, we have 
\begin{equation}~\label{eq01freqselfconsistency}
  \operatorname{freq}_1\bigl(\FD^{(a,b)}\bigr)=\frac{P}{(a+1)P+Q}.
\end{equation}

Since $P=\alpha/(\alpha+a)$
and $Q=a/(\alpha+a)$, by using $\alpha(b,a)=\alpha(a,b)=\alpha$, substituting that
into~\eqref{eq01freqselfconsistency} gives
\[
  \frac{P}{(a+1)P+Q}
  = \frac{\alpha/(\alpha+a)}{(a+1)\alpha/(\alpha+a) + a/(\alpha+a)}
  = \frac{\alpha}{(a+1)\alpha+a},
\]
which equals the target value $b/(\alpha+b)$ precisely when $ \alpha(\alpha+b) = b\bigl[(a+1)\alpha+a\bigr]$, $\alpha^2+\alpha b = ab\alpha+b\alpha+ab$ and  $\alpha^2 = ab\alpha+ab,$
which is exactly the defining equation for $\alpha$. Hence the ansatz
is self-consistent, and since~\eqref{eq01freqselfconsistency} together
with its $\sigma_b$-counterpart determines the pair of frequencies
uniquely, the
ansatz~\eqref{eq01freqcorrected} is the actual frequency vector.
\end{proof}

For example, if  $(a,b)=(1,2)$, $\alpha\approx2.7321$, and~\eqref{eq01freqcorrected}
gives $\operatorname{freq}_1(\FD^{(1,2)})\approx0.4226$. Direct computation
of $\FD^{(1,2)}_n$ for $n$ up to $9$ (length $265$) gives an empirical
frequency of letter $1$ equal to $112/265\approx0.4226$, matching to
four decimal places. For $(a,b)=(2,1)$, equation~\eqref{eq01freqcorrected}
gives $\operatorname{freq}_1(\FD^{(2,1)})=1/(1+\alpha)\approx0.2679$;
direct computation of $\FD^{(2,1)}_9$ (length $209$) gives $56/209\approx0.2679$.

\medskip

Formula~\eqref{eq01freqcorrected} reduces to the frequencies quoted
in~\eqref{eqq1biperiodicfreqs} exactly when $b=1$; for $b\ne1$ it
differs from~\eqref{eqq1biperiodicfreqs} by the presence of $b$ in
place of $1$ in the denominator. Thus, for $a\ne b$, $\operatorname{freq}_1\bigl(\FD^{(a,b)}\bigr)\ \ne\ \operatorname{freq}_1\bigl(\FD^{(b,a)}\bigr).$

\begin{lemma}~\label{lem02exponenttransform}
Let $u\in\AlphA^{*}$ be primitive, let $u_p$ be a (possibly empty) proper
prefix of $u$, let $j\geqslant 1$, and let $w=u^ju_p$. If $w$ is a factor of
$\FD^{(b,a)}$, then $\sigma_a(w)$ is a factor of $\FD^{(a,b)}$ with exponent
\begin{equation}~\label{eq01exponenttransform}
  e\bigl(\sigma_a(w)\bigr) = j + \frac{a|u_p|_0+|u_p|}{a|u|_0+|u|}.
\end{equation}
In particular, $e(\sigma_a(w)) = j+|u_p|/|u| = e(w)$ if and only if
$|u_p|_0/|u_p| = |u|_0/|u|$ whenever $u_p\ne\varepsilon$.
\end{lemma}
\begin{proof}
Since $\sigma_a$ is a morphism, $\sigma_a(w)=\sigma_a(u)^j\sigma_a(u_p)$,
and $\sigma_a(u_p)$ is a proper prefix of $\sigma_a(u)$ because $u_p$ is
a proper prefix of $u$ and $\sigma_a$ is non-erasing and
order-preserving. According to Lemma~\ref{lem01primitivity}, $\sigma_a(u)$ is
primitive, so by the standard characterisation of the smallest period
of a word of the form $v^jv_p$ with $v$ primitive and $v_p$ a proper
prefix of $v$, the smallest
period of $\sigma_a(w)$ is $|\sigma_a(u)|$, giving
\[
  e(\sigma_a(w)) = \frac{|\sigma_a(w)|}{|\sigma_a(u)|}
  = j + \frac{|\sigma_a(u_p)|}{|\sigma_a(u)|}
  = j + \frac{a|u_p|_0+|u_p|}{a|u|_0+|u|},
\]
according to Lemma~\ref{lem01length} in the last step. That $\sigma_a(w)$ is a
factor of $\FD^{(a,b)}$ follows from $w$ being a factor of $\FD^{(b,a)}$ and
$\FD^{(a,b)}=\sigma_a(\FD^{(b,a)})$. Thus, if $w$ occurs
at a block-aligned position in $\FD^{(b,a)}$, its image occupies the
corresponding block-aligned window of $\FD^{(a,b)}$.
\end{proof}

\begin{lemma}~\label{prop01celowerbound}
For $a,b \geqslant 1$, 
\[
  \CE\bigl(\FD^{(a,b)}\bigr) \ \geqslant
  \sup\Bigl\{\, j + \frac{a|u_p|_0+|u_p|}{a|u|_0+|u|}  \Bigr\},
\]
where $u^ju_p$ is a factor of $\FD^{(b,a)}$, $u$ primitive and $u_p$ is a proper prefix of $u,\ j\geqslant 1$. 
\end{lemma}
\begin{proof}
Immediate from Lemma~\ref{lem02exponenttransform} where every factor of the
stated form contributes, via $\sigma_a$, a factor of $\FD^{(a,b)}$ with
the exponent given by~\eqref{eq01exponenttransform}, and $\CE(\FD^{(a,b)})$
is the supremum of exponents over all factors of $\FD^{(a,b)}$, of
which the $\sigma_a$ images considered here are a subfamily.
\end{proof}

Actually, according to formula~\eqref{eqq1biperiodiccritical} it shows that does not reduce to the known
value $\CE(f)=2+\fphi$ when $a=b=1$
under the natural reading of $\alpha=\fphi$, giving $2+\fphi+1/\fphi=1+2\fphi$. 
Recall $s_n=\FD^{(a,b)}_{n+1}$ and $L_n=|s_n|$, where 
$L_n = d_n L_{n-1}+L_{n-2}$ with $(d_1,d_2,\ldots)=(a,b,a,b,\ldots)$.

\begin{proposition}~\label{prop01icelowerbound}
For every $n\geqslant 1$, $s_n^{d_{n+1}}s_{n-1}$ is a prefix of $\FD^{(a,b)}$, and consequently
\[
  \ice\bigl(\FD^{(a,b)}\bigr)\geqslant \sup_{n\geqslant 1}
  \left(d_{n+1}+\frac{L_{n-1}}{L_n}\right).
\]
\end{proposition}
\begin{proof}
Since $s_n$ is primitive (a standard fact for standard Sturmian words, see
\cite{Lothaire2002}), $s_{n-1}$ is a proper prefix of $s_n$, and
$s_n^{d_{n+1}}s_{n-1}=s_{n+1}$ by
Definition~\ref{def01standardsequence}, which is itself a prefix of
$\FD^{(a,b)}$. The stated exponent then follows exactly as in
Lemma~\ref{lem02exponenttransform}, by considering $u=s_n$, $u_p=s_{n-1}$,
$j=d_{n+1}$.
\end{proof}

Actually, as with $\CE$, we have verified this lower bound is not tight at
$a=b=1$; it gives $\sup_n(d_{n+1}+L_{n-1}/L_n) \to 1+1/\fphi=\fphi$, one short of the known value $\ice(f)=1+\fphi$ (see~\cite{Cassaigne2008}). Thus, the extremal prefixes realising
$\ice$ evidently extend at least one further renormalisation level
beyond $s_{n+1}$; identifying them precisely.

\medskip 

%=============================
\subsection{Duality of Return Words}~\label{sec5duality}
Having established the duality for letter frequencies and critical exponents, we now turn to return words. This subsection examines how the return-word structure of $\FD^{(a,b)}$ transforms under the substitution duality $a\leftrightarrow b$ and provides explicit descriptions of the return words for each letter. 
We begin with the established Theorem~\ref{thm01trivialinvariants} when addressing Question~\ref{q02balance}.

\begin{theorem}~\label{thm01trivialinvariants}
For every $a,b\geqslant 1$,  $\FD^{(a,b)}$ is $1$-balanced, where $B(n)\equiv1$;
and its abelian complexity is $\mathcal{AC}_{\FD^{(a,b)}}(n)\equiv2$ for
all $n\geqslant 1$.
\end{theorem}
\begin{proof}
Since $\FD^{(a,b)}$ is an aperiodic Sturmian
word for every $a,b\geqslant 1$, it is a classical theorem of
Morse--Hedlund~\cite{MorseHedlund1940} that a binary infinite word is
aperiodic and $1$-balanced if and only if it is Sturmian, which gives
$B(n)\equiv1$. For the abelian complexity, we find that $1$-balance means that for
each $n$, the set $\{|u|_1 : u\in\mathcal{L}_n(w)\}$ consists of
consecutive integers differing by at most $1$; aperiodicity of a
Sturmian word forces both values $\lfloor n\theta\rfloor+c$ and
$\lceil n\theta\rceil+c$ for the slope $\theta$ and an offset $c$ to
occur among the factors of length $n$ for every $n\ge1$, giving exactly
two abelian classes; see~\cite{RichommeSaariZamboni2010} for this
standard fact and its contrast with the Tribonacci case.
\end{proof}

Since both invariants are constant in $(a,b)$, duality under
$a\leftrightarrow b$ holds for them trivially. Thus, $B_{\FD^{(a,b)}}(n)=B_{\FD^{(b,a)}}(n)=1$ and $\mathcal{AC}_{\FD^{(a,b)}}(n)=\mathcal{AC}_{\FD^{(b,a)}}(n)=2$ for all $n$.

\begin{definition}~\label{def01returnword}
Let $w\in\AlphA^{\N}$ and let $u\in\mathcal{L}(w)$ occur at positions
$p_0<p_1<p_2<\cdots$ in $w$. The \emph{return words} of $u$ in $w$ are
the factors $w[p_i,p_{i+1})$, $i\geqslant 0$ where each such factor begins with an
occurrence of $u$ and extends to just before the next occurrence.
\end{definition}

According to Theorem~\ref{thm01vuillon}, which applies to
every $\FD^{(a,b)}$, we compute the two return words for each letter
explicitly, using the block decomposition of $\FD^{(a,b)}=\sigma_a(\FD^{(b,a)})$
from Theorem~\ref{thm01desub}.

\begin{theorem}~\label{thm01return0}
For every $a,b\geqslant 1$, the return words of $0$ in $\FD^{(a,b)}$ are exactly $\{\,0,\ 01\,\}.$
\end{theorem}
\begin{proof}
Assume that $\FD^{(a,b)}=B_1B_2B_3\cdots$ where $B_i=\sigma_a(c_i)\in
\{0^a1,\,0\}$ according to whether the $i$-th letter $c_i$ of
$\FD^{(b,a)}$ is $0$ or $1$ according to Theorem~\ref{thm01desub}. Then, we enumerate every
occurrence of $0$ and its return word:

\medskip 

\noindent \textbf{Case 1.} Within a block $B_i=0^a1$ ($c_i=0$). In this case, the $a$ occurrences of $0$
in this block are consecutive positions, so the return word from the
$k$-th to the $(k{+}1)$-th of these ($1\le k\le a-1$) is simply $0$.

\medskip 

\noindent \textbf{Case 2.} From the last $0$ of a block $B_i=0^a1$. The next character is
the terminal $1$ of $B_i$, and the next occurrence of $0$ is the first
character of $B_{i+1}$ (whatever $B_{i+1}$ is, since both possible
blocks $0^a1$ and $0$ begin with $0$). Hence the return word is exactly
$01$, independently of $B_{i+1}$.

\medskip 


\noindent \textbf{Case 3.} From a block $B_i=0$ ($c_i=1$). This is a single occurrence of
$0$ with nothing else in the block; the next character is the first
character of $B_{i+1}$, which is again $0$ regardless of which block
$B_{i+1}$ is. Hence the return word is exactly $0$.

Every occurrence of $0$ in $\FD^{(a,b)}$ falls into exactly one of these
three cases, all of which give return word $0$ or $01$, and both occur. This proves the claim for
$a\geqslant2$.
\end{proof}
For $a=1$, blocks $0^a1=01$ contain no internal repetition,
but the boundary case alone already produces both $0$, from a $B_i=0$
block and $01$ from the boundary after any $0^11$ block, so both
return words still occur.

\begin{theorem}~\label{thm02return1}
For every $a,b\geqslant 1$, the return words of $1$ in $\FD^{(a,b)}$ are exactly $\{\,1\,0^a,\ 1\,0^{a+1}\,\}.$
\end{theorem}
\begin{proof}
Assume that every occurrence of $1$ in $\FD^{(a,b)}$ occurs as the terminal letter of
a block $B_i=0^a1$, where it corresponds to an occurrence of the letter
$0$ at position $i$ in $\FD^{(b,a)}$. By Theorem~\ref{thm01return0}
applied to $\FD^{(b,a)}$, the return words of $0$ in $\FD^{(b,a)}$ are $0$
and $01$. Then, between two consecutive occurrences of $0$ in
$\FD^{(b,a)}$ there are either zero or exactly one intervening $1$.

\medskip 


\noindent \textbf{Case 1.} If consecutive source $0$'s are
adjacent (return word $0$ in $\FD^{(b,a)}$), the corresponding blocks
$B_i=0^a1$ and $B_{i+1}=0^a1$ are themselves consecutive, so $\FD^{(a,b)}$ from one terminal $1$ to the next consists exactly of
the $a$ leading zeros of $B_{i+1}$ followed by its terminal $1$. Thus, the
return word, starting at the occurrence of
$1$, is $1\cdot0^a$.

\medskip 


\noindent \textbf{Case 2.} If there is exactly one intervening source $1$ (return word $01$ in $\FD^{(b,a)}$), the corresponding image contains one extra single-letter
block $0$ from that source $1$ inserted before the next $0^a1$ block;
the gap becomes $0$ from the intervening block together with the $a$
leading zeros of the next $0^a1$ block, yielding $a+1$ zeros before the
next $1$. Thus, the return word is $1\cdot0^{a+1}$.

Since both return words of Case 1 and Case 2 of $0$ occur in $\FD^{(b,a)}$, both $1\cdot0^a$ and $1\cdot0^{a+1}$
occur as return words of $1$ in $\FD^{(a,b)}$, and by
Theorem~\ref{thm01vuillon} there are no others.
\end{proof}

Assume that $a=b=1$, then acccording to Theorem~\ref{thm01return0} predicts return words of $1$
in the classical Fibonacci word $f$ equal to $\{10,100\}$. Directly
from $f=0100101001001010\cdots$, the occurrences of $1$ are at
positions $1,4,6,9,12,14,\ldots$, giving return words $f[1,4)=100$,
$f[4,6)=10$, $f[6,9)=100$, $f[9,12)=100$, $f[12,14)=10,\ldots$,
confirming $\{10,100\}$.


\begin{corollary}~\label{cor01returnduality}
The return words of $1$ in $\FD^{(a,b)}$ and $\FD^{(b,a)}$ are $\{1\,0^a,\ 1\,0^{a+1}\}$ and 
$\{1\,0^b,\ 1\,0^{b+1}\}$ respectively. Exchanging $a$ and $b$ acts on the return-word set of $1$ by exactly replacing the exponent $a$ by $b$, while the
return-word set of $0$, $\{0,01\}$, remains unchanged. 
\end{corollary}

Corollary~\ref{cor01returnduality} shows that the letter $0$ plays a
completely symmetric role: its return words never depend on $a$ or $b$,
while the letter $1$ carries the entire $(a,b)$-dependence through the
first parameter of the word under consideration. This constitutes a
second, independent instance of the asymmetric letter roles under $\sigma_a$.

The lengths
$a+1$ and $a+2$ of these return words are also exactly the two
consecutive integers between which the length ratios $L_n/L_{n-1}$ of
Section~\ref{sec01sadic} oscillate. The decomposition of
Theorem~\ref{thm01return0} is the length-$1$ level of the
\emph{Ostrowski numeration system} associated with the continued
fraction $\chi(a,b)=[a,b,a,b,\ldots]$ of
Proposition~\ref{prop01AlphAacf}, with higher return words of the
factors $s_n$ corresponding to higher digits. 

%=================================
\subsection{Duality of Palindromic Structure}~\label{sec01palduality}
This subsection addresses the duality of the palindromic structure. We connect the palindromicity results for $\FD^{(a,b)}$ to the classical theory of central words and discuss consequences for palindromic complexity.

\begin{theorem}~\label{thm01centralpalindrome}
For every $a,b\geqslant 1$ and every $n\geqslant 2$, $\bigl(\FD^{(a,b)}_n\bigr)^{-}$ is
a palindrome.
\end{theorem}
\begin{proof}
\noindent \textbf{Method 1.} 
By Proposition~\ref{prop01fabisstandard}, the sequence $(\FD^{(a,b)}_n)_{n\ge0}$ coincides with the standard sequence $(s_n)_{n\ge0}$ associated with the directive sequence $(d_1,d_2,\ldots)=(a,b,a,b,\ldots)$. In particular, for $n\ge 2$ we have $\FD^{(a,b)}_n=s_{n-1}$.

It is a classical result in the theory of Sturmian words that every term $s_k$ ($k\ge 1$) of the standard sequence satisfies the following property: the word obtained by removing its last two letters, $(s_k)^- \ = \ s_k[0..|s_k|-3],$
is a central palindrome. This fact follows from the recursive construction of standard Sturmian words and the palindromic closure properties of Christoffel words / standard words (see, e.g., Lothaire~\cite{Lothaire2002}, Ch.~2). Since $n\geqslant 2$ implies $n-1\geqslant  1$, the word $(\FD^{(a,b)}_n)^-$ is a central palindrome.

\medskip

\noindent \textbf{Method 2.}  Since $s_{k-2}$ is a prefix of $s_{k-1}$, assume that  $s_k = t_k\,xy$
for its last two letters $x,y$.

\medskip

\noindent \textbf{Claim 1.} For $k\geqslant 1$, the last two letters of $s_k$ are $(0,1)$ if $k$ is odd, and $(1,0)$ if $k$ is even.

\noindent \textit{Proof.} 
Since $s_1 = 0^{d_1}1$, and $s_1$ ends with $01$. Then, $s_2 = s_1^{d_2}s_0.$
The final block is $s_0=0$, preceded by the last letter of $s_1$ which is $1$. Thus, $s_2$ ends with $10$. For $k\geqslant 3$, we have $s_k = s_{k-1}^{d_k}s_{k-2}.$

Since $|s_{k-2}|\ge |s_1|\geqslant 2$, the last two letters of $s_k$ are exactly the last two letters of the final block $s_{k-2}$. By strong induction, the claim for $k-2$ follows, with the same parity as $k$. \hfill \qedsymbol

\medskip


\noindent \textbf{Claim 2.} For every $k\geqslant 1$, the words $s_{k-1}s_{k-2}$ and $s_{k-2}s_{k-1}$ have the same length and the same degree in every letter except the last two, which are exchanged, if $s_{k-1}s_{k-2}=R_k\,xy$. Then, $s_{k-2}s_{k-1}=R_k\,yx.$

\noindent \textit{Proof.} 
For $k=1$, we have $s_0s_{-1}=01$ and $s_{-1}s_0=10,$ so the claim is immediate: the two words agree on the empty prefix and only the last two letters are swapped. Assume $k\geqslant 2$. By using $s_{k-1}=s_{k-2}^{d_{k-1}}s_{k-3},$
we obtain  $s_{k-1}s_{k-2}=s_{k-2}^{d_{k-1}-1}\cdot (s_{k-2}s_{k-3}).$
Similarly, $s_{k-2}s_{k-1}=s_{k-2}^{d_{k-1}-1}\cdot (s_{k-2}s_{k-3}).$ 

Thus both words are the same prefix $s_{k-2}^{d_{k-1}-1}$ followed by the pair $(s_{k-2}s_{k-3})$, which is exactly the claim at level $k-1$. By induction, that pair agrees except in its last two letters, which are swapped; prepending the identical block preserves this.
\hfill \qedsymbol 

Thus, according to both claims $\bigl(\FD^{(a,b)}_n\bigr)^{-}$ is
a palindrome.
\end{proof}

This is an instance of the classical fact that the levels of a standard
Sturmian sequence, the same statement as the palindromicity of
$\Phi(\FD^{(a,b)}_n)$ established in~\cite{Ramirez2015}, where 
$\Phi(\FD^{(a,b)}_n)=\bigl(\FD^{(a,b)}_n\bigr)^{-}$ or a word differing
from it by a bounded, easily tracked correction.  

\begin{corollary}~\label{cor002arbitrary}
$\FD^{(a,b)}$ has palindromic prefixes of arbitrarily large length; in
particular $P_{\FD^{(a,b)}}(|\FD^{(a,b)}_n|-2)\geqslant 1$ for every $n\geqslant 2$.
\end{corollary}
\begin{proof}
Since $\FD^{(a,b)}_n$ is a prefix of $\FD^{(a,b)}$ by Proposition~\ref{prop01fabisstandard} and the discussion preceding
it, $\bigl(\FD^{(a,b)}_n\bigr)^{-}$ is a palindromic factor of $\FD^{(a,b)}$. As $|\FD^{(a,b)}_n|\to\infty$, the claim follows.
\end{proof}

Corollary~\ref{cor002arbitrary} is consistent with,
and gives an independent confirmation of, the characteristic standard
status of $\FD^{(a,b)}$ already established in Corollary~\ref{cor01sturmian}.

%==============================
%========================
\section{The Continued Fraction Expansion of the Slope of \texorpdfstring{$\FD^{(a,b)}$}{FD(a,b)}}\label{sec01continuedfraction}

This section determines the continued fraction expansion of the slope $\theta^{(a,b)}$ of the biperiodic word $\FD^{(a,b)}$. The explicit expansion allows us to derive several quantitative properties of these words independently and reveals that both the slope and the index $\operatorname{Ind}(\FD^{(a,b)})$ depend on the pair $(a,b)$ only through the product $ab$.

We determine the continued fraction expansion of the
slope of the biperiodic Fibonacci word $\FD^{(a,b)}$, i.e.\ of the
irrational number
\begin{equation}~\label{eqq01slopedef}
  \theta^{(a,b)}\;:=\; \lim_{n\to\infty}\frac{|\FD^{(a,b)}_n|_1}{|\FD^{(a,b)}_n|}
  \;=\; \frac{1}{1+\AlphA },
\end{equation}
where $\alpha=\alpha(a,b)$ is as in \eqref{eqq1AlphAadef}. This
result is the technical core of the paper: it lets us re-derive
$\operatorname{CE}(\FD^{(a,b)})$ independently,
rather than treating~\eqref{eqq1biperiodiccritical} as an unverified
black box.

Throughout, $a,b$ are positive integers, and for a real number
$x>1$ we write $x=[a_0;a_1,a_2,\dots]$ for its (simple) continued
fraction expansion, and $x=[\overline{a_1,\dots,a_k}]$ if the
expansion is purely periodic with period $(a_1,\dots,a_k)$. 
Recall that a quadratic irrational $x>1$ has a \emph{purely periodic}
continued fraction expansion if and only if $x$ is \emph{reduced},
i.e.\ its conjugate $x'$ satisfies $-1<x'<0$.

\begin{proposition}~\label{prop01alphareduced}
For all integers $a,b\geqslant 1$, $\AlphA =\AlphA (a,b)$ is a reduced
quadratic irrational.
\end{proposition}
\begin{proof}
Since $a,b\geqslant 1$ we have $ab\geqslant 1$, so $a^2b^2+4ab>a^2b^2$. Hence, 
$\AlphA =\tfrac{ab+\sqrt{a^2b^2+4ab}}{2} > ab \geqslant 1$. For the conjugate, since $4ab>0$ we have
$\sqrt{a^2b^2+4ab}>ab$. Then, 
\[
  \AlphA' = \frac{ab-\sqrt{a^2b^2+4ab}}{2} < 0 .
\]
It remains to show $\AlphA'>-1$. Note that 
$\sqrt{a^2b^2+4ab} < ab+2$. As both sides are positive, this is
equivalent to
\[
  a^2b^2+4ab \;<\; (ab+2)^2 \;=\; a^2b^2+4ab+4,
\]
which holds for every $a,b$. Hence $-1<\AlphA '<0$, so $\AlphA $ is
reduced.
\end{proof}

\begin{lemma}~\label{lem01alphaquadratic}
The constant $\AlphA =\AlphA (a,b)=(\sqrt{a^2b^2+4ab}+ab)/2$ is
the unique positive root of
\begin{equation}~\label{eqq01alphaminpoly}
  t^2 - ab\,t - ab = 0.
\end{equation}
Consequently its algebraic conjugate is
$\AlphA ' = (ab-\sqrt{a^2b^2+4ab})/2$, and
\[
  \AlphA +\AlphA ' = ab, \qquad \AlphA \,\AlphA ' = -ab.
\]
\end{lemma}
\begin{proof}
Immediate from the quadratic formula applied to
$t^2-ab\,t-ab=0$, whose discriminant is $a^2b^2+4ab$; the positive
root is exactly the stated expression for $\AlphA $.
\end{proof}

Lemma~\ref{lem01alphaquadratic} shows that $\AlphA (a,b)$ depends on
$a$ and $b$ only through the product $ab$: if $ab=a'b'$ then
$\AlphA (a,b)=\AlphA (a',b')$, even though $(a,b)\neq(a',b')$. For
$a=b=1$ this recovers $\AlphA =\varphi$, the golden ratio, as the
positive root of $t^2-t-1=0$.

\begin{theorem}~\label{thm01alphacf}
For all integers $a,b\geqslant 1$,
\begin{equation}~\label{eq01alphacf}
  \AlphA (a,b) \;=\; [\overline{ab,\,1}].
\end{equation}
\end{theorem}
\begin{proof}
By Proposition~\ref{prop01alphareduced}, $\AlphA $ has a unique
purely periodic continued fraction expansion. It therefore suffices
to check that $x:=[\overline{ab,1}]$ satisfies the same minimal
polynomial~\eqref{eqq01alphaminpoly} as $\AlphA $, and that $x>1$ with
conjugate in $(-1,0)$.

For $x\geqslant 1$, 
\[
  x = ab + \cfrac{1}{1+\cfrac{1}{x}} = ab+\frac{x}{x+1}.
\]
Clearing the denominator yields $x(x+1) = ab(x+1) + x$, so 
\[
  x^2 + x = abx + ab + x \qquad \Longrightarrow \qquad x^2 - ab\,x - ab = 0,
\]
which is exactly~\eqref{eqq01alphaminpoly}. Since $x=[\overline{ab,1}]>ab\ge1$
and the two roots of~\eqref{eqq01alphaminpoly} are $\AlphA >1$ and
$\AlphA '\in(-1,0)$, it follows that $x=\AlphA $.
\end{proof}

\begin{example}
According to Theorem~\ref{thm01alphacf}, 
\begin{enumerate}
    \item For $a=b=1$ (so $ab=1$), we have $\AlphA =[\overline{1,1}]=[\overline{1}]=\varphi$,
recovering the classical golden ratio.
    \item For $a=2,b=3$ (so $ab=6$), we have $\AlphA =[\overline{6,1}]$. Numerically,
$\AlphA \approx 6.8730$, which agrees with
$\AlphA (2,3)=\bigl(\sqrt{36+24}+6\bigr)/2\approx 6.8730$.
\end{enumerate}
\end{example}

We now pass from $\AlphA $ to the actual slope $\theta^{(a,b)}=1/(1+\AlphA )$
of $\FD^{(a,b)}$.

\begin{theorem}~\label{thm01slopecf}
For all integers $a,b\geqslant 1$, the slope of the biperiodic Fibonacci
word $\FD^{(a,b)}$ has continued fraction expansion
\begin{equation}~\label{eqq2slopecf}
  \theta^{(a,b)} \;=\; \bigl[\,0;\; ab+1,\; \overline{1,\,ab}\,\bigr].
\end{equation}
\end{theorem}
\begin{proof}
According to Theorem~\ref{thm01alphacf}, $\AlphA =[\overline{ab,1}]$, so 
$\lfloor\AlphA \rfloor=ab$ and the fractional part is
\[
  \{\AlphA \} \;=\; \frac{1}{[\,\overline{1,\,ab}\,]}\,.
\]
Thus,
\[
  1+\AlphA  = (ab+1) + \{\AlphA \} = \bigl[ab+1;\,\overline{1,\,ab}\,\bigr].
\]
Since $1+\AlphA >1$, the reciprocal rule for continued fractions gives
\[
  \theta^{(a,b)} = \frac{1}{1+\AlphA }
  = \bigl[\,0;\, ab+1,\,\overline{1,\,ab}\,\bigr].
\]
\end{proof}

\begin{example}
For $a=b=1$, Theorem~\ref{thm01slopecf} gives
$\theta^{(1,1)} = [0;2,\overline{1,1}] = [0;2,\overline{1}]$, which is
the continued fraction expansion of $1/\varphi^2\approx0.381966$,
the known slope of the classical Fibonacci word. For $a=2,b=3$, we obtain $\theta^{(2,3)}=[0;7,\overline{1,6}]$.
\end{example}

The slope $\theta^{(a,b)}$, and
hence the continued-fraction data that determines it, depends on the
pair $(a,b)$ only through the integer $ab$. In particular, whenever
$ab=a'b'$ for two distinct pairs $(a,b)\neq(a',b')$ (for example, 
$(a,b)=(2,3)$ and $(a',b')=(1,6)$), the words
$\FD^{(a,b)}$ and $\FD^{(a',b')}$ have the same slope
$\theta^{(a,b)}=\theta^{(a',b')}$ and therefore the same critical exponent. 
For integers $a,b\geqslant 1$, the finite biperiodic Fibonacci words
$\FD^{(a,b,n)}$ are defined by
\[
  \FD^{(a,b,0)}=\varepsilon,\quad \FD^{(a,b,1)}=0,\quad \FD^{(a,b,2)}=0^{a-1}1,
\]
\[
  \FD^{(a,b,n)} =
  \begin{cases}
    \bigl(\FD^{(a,b,n-1)}\bigr)^{a} \FD^{(a,b,n-2)}, & n\geqslant 3 \text{ even},\\[2pt]
    \bigl(\FD^{(a,b,n-1)}\bigr)^{b} \FD^{(a,b,n-2)}, & n\geqslant 3 \text{ odd},
  \end{cases}
\]
and $\FD^{(a,b)}:=\lim_{n\to\infty}\FD^{(a,b,n)}$. Their length
$q_n:=|\FD^{(a,b,n)}|$ satisfies the recurrence
\begin{equation}~\label{eqq01qnrecurrence}
  q_0=0,\ q_1=1,\qquad
  q_n = \begin{cases} a\,q_{n-1}+q_{n-2}, & n\geqslant 2\text{ even},\\
                       b\,q_{n-1}+q_{n-2}, & n\geqslant 2\text{ odd}.\end{cases}
\end{equation}

\begin{proposition}~\label{prop001convergentidentity}
For all $n\geqslant 0$, the length $q_n=|\FD^{(a,b,n)}|$ defined
by~\eqref{eqq01qnrecurrence} coincides with the denominator $q_{n-1}$ of the $(n-1)$-th convergent to the continued fraction $\zeta=[0;\overline{a,b}]$ (with the convention $q_{-1}=0$).
\end{proposition}
\begin{proof}
The sequences satisfy the same initial conditions and the same parity-dependent two-term linear recurrence.
\end{proof}

Now let $M:=\max(a,b)$ and $P:=ab$. Define
\begin{equation}~\label{eqq01indMP}
  \operatorname{Ind}\bigl(\FD^{(a,b)}\bigr) = 2+M\left(1+\frac1{\AlphA (P)}\right),
  \qquad   \AlphA (P)=\frac{P+\sqrt{P^2+4P}}{2}.
\end{equation}

\begin{lemma}~\label{lem01alphamonotone}
The function $\AlphA (P)$ is strictly increasing for $P>0$.
\end{lemma}
\begin{proof}
Differentiating gives $\AlphA '(P) = \frac12\left(1+\frac{P+2}{\sqrt{P^2+4P}}\right) >0$
for all $P>0$.
\end{proof}

\begin{lemma}~\label{lem01lockstructure}
Let $M\geqslant 1$ and let $(a,b)$ range over all pairs with $\max(a,b)=M$. Then
$\operatorname{Ind}(\FD^{(a,b)})$ is a strictly decreasing function of $m=\min(a,b)$.
\end{lemma}
\begin{proof}
For fixed $M$, $\operatorname{Ind}(\FD^{(a,b)})=2+M(1+1/\AlphA (Mm))$ depends on $m$ only through $\AlphA (Mm)$. Since $\AlphA $ is strictly increasing (Lemma~\ref{lem01alphamonotone}), $\operatorname{Ind}$ is strictly decreasing in $m$.
\end{proof}

\begin{lemma}~\label{lem03diagonalincreasing}
The function $g(M):=\operatorname{Ind}(\FD^{(M,M)})$ is strictly increasing for $M>0$. Explicitly,
\begin{equation}~\label{eqq1lem03diagonalincreasing}
g(M) = 2+M+\dfrac{2}{M+\sqrt{M^2+4}}.
\end{equation}
\end{lemma}
\begin{proof}
On the diagonal, $P=M^2$ and
\[
  \frac{M}{\AlphA (M^2)} = \frac{2}{M+\sqrt{M^2+4}}.
\]
Let $s(M)=M+\sqrt{M^2+4}$. Then $g(M)=2+M+2/s(M)$. Differentiating and verifying $g'(M)>0$ for $M>0$ proceeds exactly as in the original calculation.
\end{proof}
We illustrate that by Algorithm~\ref{alg01ind} for $\operatorname{Ind}(\FD^{(a,b)})$ based on~\eqref{eqq01indMP}.  To reduces the critical exponent to three elementary arithmetic operations once $a,b$ are fixed, which makes it convenient both for generating tables and for the numerical verification.
\begin{algorithm}[H]
\small
\caption{Compute $\operatorname{Ind}(\FD^{(a,b)})$}~\label{alg01ind}
\begin{algorithmic}[1]
\Require positive integers $a,b$
\Ensure $\operatorname{Ind}(f^{(a,b)})$
\State $P \gets a\cdot b$
\State $M \gets \max(a,b)$
\State $\AlphA \gets \dfrac{P+\sqrt{P^2+4P}}{2}$
\State \Return $2 + M\left(1+\dfrac{1}{\AlphA}\right)$
\end{algorithmic}
\end{algorithm}

\begin{theorem}~\label{thm01globalmin}
Among all $\FD^{(a,b)}$ with $a,b\geqslant 1$ integers,
$\operatorname{Ind}(\FD^{(a,b)})$ attains its global minimum uniquely at
$(a,b)=(1,1)$, with value $2+\varphi$.
\end{theorem}
\begin{proof}
Let $M=\max(a,b)$. By Lemma~\ref{lem01lockstructure},
$\operatorname{Ind}(\FD^{(a,b)}) \geqslant g(M)$, with equality if and only if $a=b=M$. By Lemma~\ref{lem03diagonalincreasing}, $g(M)\geqslant g(1)$ with equality if and only if $M=1$. Hence $\operatorname{Ind}(\FD^{(a,b)})\geqslant g(1)=2+\varphi$, with equality if and only if $(a,b)=(1,1)$.
\end{proof}

\begin{corollary}~\label{cor01symmetry}
For $a,b\geqslant 1$, $\operatorname{Ind}(\FD^{(a,b)})=\operatorname{Ind}(\FD^{(b,a)})$.
\end{corollary}
\begin{proof}
Immediate from~\eqref{eqq01indMP}, since both $M=\max(a,b)$ and $P=ab$ are symmetric in $(a,b)$.
\end{proof}


To generate the underlying word itself — needed, for instance, to check
Algorithm~\ref{alg01ind}'s output against an empirical repetition search. Algorithm~\ref{alg01wordgen}
builds $\FD^{(a,b,n)}$ bottom-up from the two base cases, by applying the
parity-dependent exponentiation rule at each step; its output length
$|w_n|$ coincides with $q_n$ from~\eqref{eqq01qnrecurrence} by
construction.
\begin{algorithm}[H]
\small
\caption{Generate the finite biperiodic Fibonacci word $\FD^{(a,b,n)}$}
\label{alg01wordgen}
\begin{algorithmic}[1]
\Require positive integers $a,b$; target index $n$
\Ensure the word $\FD^{(a,b,n)}$ as a string
\State $w_0 \gets \varepsilon$,\quad $w_1 \gets \texttt{"0"}$,\quad
$w_2 \gets \texttt{"0"}\times(a-1) \,\Vert\, \texttt{"1"}$
\For{$k = 3$ \textbf{to} $n$}
  \If{$k$ is even}
    \State $w_k \gets w_{k-1}^{\,a} \,\Vert\, w_{k-2}$
  \Else
    \State $w_k \gets w_{k-1}^{\,b} \,\Vert\, w_{k-2}$
  \EndIf
\EndFor
\State \Return $w_n$
\end{algorithmic}
\end{algorithm}
%%%%%%%%%%%%%%%%%%%%%%%%%%%%
\section{Conclusion and Discussion}

We set out from a single observation: the letter-frequency parameter
$\alpha(a,b)$ of the biperiodic Fibonacci word $\FD^{(a,b)}$ is symmetric
in $a$ and $b$, while the quoted critical exponent formula is not. This
suggested some structural correspondence between $\FD^{(a,b)}$ and
$\FD^{(b,a)}$ that a symmetric algebraic quantity alone could not explain. 
Realising $\FD^{(a,b)}$ as the standard Sturmian sequence associated with the
purely periodic continued fraction $[a,b,a,b,\ldots]$ shows that $\FD^{(a,b)} = \sigma_a\bigl(\FD^{(b,a)}\bigr)$
exactly, with no relabelling or bounded correction term
(Theorem~\ref{thm01desub}), where $\sigma_a\colon 0\mapsto 0^a1,\,1\mapsto 0$.
Corollary~\ref{cor01sturmian} emphasizes that $\FD^{(a,b)}$ is Sturmian for
\emph{every} $a,b\geqslant 1$, so that balance and abelian complexity are
constantly equal to $1$ and $2$, respectively. 

Return words admit a complete, explicitly verified description
Theorems~\ref{thm01return0} and \ref{thm02return1}, those of $0$ are
$\{0,01\}$ independently of $a$ and $b$, while those of $1$ are
$\{1\cdot0^a,\,1\cdot0^{a+1}\}$. Thus, the exchange $a\leftrightarrow b$ acts on
them by literally substituting the exponent $a\mapsto b$
(Corollary~\ref{cor01returnduality}), providing the cleanest and most
explicit instance of the paper's central mechanism. Palindromic prefixes of arbitrarily large length exist
(Theorem~\ref{thm01centralpalindrome}), which is consistent with, and gives
an independent proof of, the standardness established in
Corollary~\ref{cor01sturmian}.

The transformation of every combinatorial invariant of $\FD^{(a,b)}$ under
$a\leftrightarrow b$ is governed by one explicit morphism, $\sigma_a$,
and the asymmetry of that transformation traces back to the same
length-redistribution mechanism in every case we examined. This provides,
we believe, a more satisfying explanation than a case-by-case analysis of
individual invariants.

\subsection{Open problems}

While this paper establishes several duality results for the family of biperiodic Fibonacci words, a number of natural and important questions remain open.

\begin{openproblem}~\label{op03palcomplexity}
Compute the palindromic complexity $P_{\FD^{(a,b)}}(n)$ for every $n$, beyond the special lengths $|\FD^{(a,b)}_n|-2$ treated in Corollary~\ref{cor002arbitrary}. 
For the classical case $a=b=1$, computation on the Fibonacci word yields $P(1)=2$, $P(2)=1$, $P(3)=2$, $P(4)=1$, suggesting the alternating pattern
\[
  P(n) = 
  \begin{cases}
    2 & \text{if $n$ is odd}, \\
    1 & \text{if $n$ is even, $n\geqslant 2$}.
  \end{cases}
\]
It is open whether an analogous closed-form expression exists for general $(a,b)$ and how it depends on the parameters.
\end{openproblem}

Any duality result for palindromic complexity under the exchange $a\leftrightarrow b$ must account for the parity-dependent endings ($\cdots 01$ or $\cdots 10$) of the standard prefixes $\FD^{(a,b)}_n$. This appears to require a finer analysis than the block decomposition that sufficed for frequencies and return words.

In addition, the following technical questions are left unresolved:

\begin{enumerate}
\item \textbf{Critical exponents.} 
  The lower bounds obtained in Propositions~\ref{prop01celowerbound} and~\ref{prop01icelowerbound} for $\CE(\FD^{(a,b)})$ and $\ice(\FD^{(a,b)})$ are not tight. Identifying the exact extremal repetitions and deriving closed-form expressions remains a central open problem.

\item \textbf{Palindromic complexity and duality.} 
  A complete closed-form description of $P_{\FD^{(a,b)}}(n)$ for all $n$, together with its behaviour under $a\leftrightarrow b$, as posed in Open Problem~\ref{op03palcomplexity}.

\item \textbf{Identification with the palindromisation map.} 
  We conjecture that the palindromisation operator $\Phi$ used in~\cite{Ramirez2015} coincides, up to a bounded correction, with the deletion of the last two letters of $\FD^{(a,b)}_n$. Direct verification of this identification is desirable.
\end{enumerate}

\subsection{Directions for future work}

Beyond resolving the items above, three broader directions follow
naturally from the framework developed here.

The return-word lengths $a+1$ and $a+2$ computed in
Section~\ref{sec5duality} are the base level of the Ostrowski
numeration system attached to $\chi(a,b)=[a,b,a,b,\ldots]$.
Developing this numeration system in full (i.e., return words of the
higher-level factors $s_n$, not just of single letters) would likely resolve
the exact critical exponent and palindromic complexity problems simultaneously, since both are
classically controlled by the same numeration data for other Sturmian
families.

The desubstitution mechanism of
Theorem~\ref{thm01desub} relies only on the fact that the directive sequence of
$\FD^{(a,b)}$ is periodic with period $2$. The same argument applies
verbatim to any $k$-periodic directive sequence
$(d_1,\ldots,d_k,d_1,\ldots,d_k,\ldots)$. Desubstitution by
$\sigma_{d_1}$ yields the word with cyclically shifted directive sequence.
This suggests a full \emph{cyclic duality} of order $k$ acting on such families, of which the present results are the $k=2$ case. Whether the return-word description of
Section~\ref{sec5duality} generalises cleanly for $k>2$ remains, to our knowledge, unexplored.

Finally, Ram\'irez, and Rubiano, original motivation for $\FD^{(a,b)}$ was geometric -- an
associated fractal curve and, in a later extension, a family of
double-square polyominoes~\cite{FibSnowflake2024}. It would be natural
to ask whether the duality map $\sigma_a$ has a geometric counterpart
on the fractal curves themselves.

The puzzle we started with turned out to be based on an incorrect
premise: letter frequency is not invariant under
$a\leftrightarrow b$. We regard the resolution as a better outcome rather than explaining one isolated coincidence, the
Desubstitution Theorem explains, uniformly and exactly, why every
combinatorial statistic of $\FD^{(a,b)}$ that depends on the interaction
between letter identity and block length fails to be symmetric under
parameter exchange, and it does so through a single, explicit,
computable morphism.

%%%%%%%%%%%%%%%%%%%%%%%%%%%%%%%
%%%%%%%%%%%%%%%%%%%%%%%%%%%%%%%%%%


%%%%%%%%%%%%%%55

\begin{thebibliography}{99}

\bibitem{Abdullah2023Meisami}
D.~Abdullah and M.~Meisami,
``Factor and arithmetic complexity of concatenating the n!'',
\emph{Chebyshevskii Sbornik} \textbf{24}(4), 341--344 (2023).
\doi{10.22405/2226-8383-2023-24-4-341-344}

\bibitem{Abdullah2025Hamoud}
D.~Abdullah and J.~Hamoud,
``Generalized Natural Density DF (Fk) of Fibonacci Word'',
\emph{Vestnik KRAUNC. Fiziko-Matematicheskie Nauki} \textbf{52}(3), 7--23 (2025).

\bibitem{AllouecheShallit2003}
J.-P.~Allouche and J.~Shallit,
\emph{Automatic Sequences: Theory, Applications, Generalizations},
Cambridge University Press, 2003.

\bibitem{BalkovaBrindaTurek2011}
L.~Balkov\'a, K.~B\v{r}inda, and O.~Turek,
``Abelian complexity of infinite words associated with quadratic Parry numbers'',
\emph{Theoretical Computer Science} \textbf{412}, 6252--6260 (2011).

\bibitem{FibSnowflake2024}
D.~Bortz, N.~Cummings, S.~Gao, E.~Jaffe, L.~Mai, B.~Steinhurst, and P.~Tillotson,
``Self-Similar Structure of k-and Biperiodic Fibonacci Words'',
\emph{PUMP Journal of Undergraduate Research} \textbf{7}, 1--13 (2024).

\bibitem{Cassaigne2008}
J.~Cassaigne,
``On extremal properties of the Fibonacci word'',
\emph{RAIRO -- Theoretical Informatics and Applications} \textbf{42}(4), 701--715 (2008).

\bibitem{CurrieRampersad2008}
J.~Currie and N.~Rampersad,
``For each $\alpha>2$ there is an infinite binary word with critical exponent $\alpha$'',
\emph{Electronic Journal of Combinatorics} \textbf{15}, \#R119 (2008).

\bibitem{DamanikLenz2002}
D.~Damanik and D.~Lenz,
``The index of Sturmian sequences'',
\emph{Journal of Combinatorial Theory, Series A} \textbf{100}(1), 189--195 (2002).

\bibitem{Dumaine2009}
A.~M.~Dumaine,
``The Fibonacci Word fractal'', hal-00367972, 2009.

\bibitem{EdsonYayenie2009}
M.~Edson and O.~Yayenie,
``A new generalization of Fibonacci sequence and extended Binet's formula'',
\emph{Integers} \textbf{9}(6), 639--654 (2009).

\bibitem{Lothaire2002}
M.~Lothaire,
\emph{Algebraic Combinatorics on Words},
Encyclopedia of Mathematics and its Applications, vol.~90,
Cambridge University Press, 2002.

\bibitem{MignosiPirillo1992}
F.~Mignosi and G.~Pirillo,
``Repetitions in the Fibonacci infinite word'',
\emph{RAIRO -- Informatique Th\'eorique et Applications} \textbf{26}(3), 199--204 (1992).

\bibitem{MorseHedlund1940}
M.~Morse and G.~A.~Hedlund,
``Symbolic dynamics II: Sturmian trajectories'',
\emph{American Journal of Mathematics} \textbf{62}, 1--42 (1940).

\bibitem{MousaviShallit2015}
H.~Mousavi and J.~Shallit,
``Mechanical proofs of properties of the Fibonacci word'',
in \emph{Combinatorics on Words (WORDS 2015)},
Lecture Notes in Computer Science, vol.~9304, Springer, pp.~170--190, 2015.

\bibitem{Ramirez2015}
J.~L.~Ram\'irez and G.~N.~Rubiano,
``Biperiodic Fibonacci Word and its Fractal Curve'',
\emph{Acta Polytechnica} \textbf{55}(1), 50--58 (2015).
\doi{10.14311/ap.2015.55.0050}

\bibitem{RichommeSaariZamboni2010}
G.~Richomme, K.~Saari, and L.~Q.~Zamboni,
``Balance and abelian complexity of the Tribonacci word'',
\emph{Advances in Applied Mathematics} \textbf{45}(2), 212--231 (2010).

\bibitem{Vuillon2001}
L.~Vuillon,
``A characterization of Sturmian words by return words'',
\emph{European Journal of Combinatorics} \textbf{22}(2), 263--275 (2001).

\end{thebibliography}










 
\end{document}
